\documentclass[10pt,reqno]{amsart}

\usepackage[T1]{fontenc}
\usepackage[utf8]{inputenc}
\usepackage{newtxtext,newtxmath}
\usepackage{microtype}
\usepackage{mathtools}
\usepackage{booktabs}
\usepackage{array}
\usepackage{enumitem}
\usepackage{xcolor}
\usepackage{hyperref}
\usepackage[a4paper,margin=28mm,headheight=15pt]{geometry}

\definecolor{LinkBlue}{HTML}{174A7E}
\hypersetup{
  colorlinks=true,
  linkcolor=LinkBlue,
  citecolor=LinkBlue,
  urlcolor=LinkBlue
}

\setlength{\parskip}{2pt}
\setlist[itemize]{leftmargin=6mm,itemsep=1.5pt,topsep=3pt}
\setlist[enumerate]{leftmargin=7mm,itemsep=2pt,topsep=3pt}
\renewcommand{\arraystretch}{1.15}
\emergencystretch=2em

\newtheorem{theorem}{Theorem}[section]
\newtheorem{proposition}[theorem]{Proposition}
\newtheorem{lemma}[theorem]{Lemma}
\newtheorem{corollary}[theorem]{Corollary}
\theoremstyle{definition}
\newtheorem{definition}[theorem]{Definition}
\newtheorem{remark}[theorem]{Remark}

\title[An Irrational Non-$\mathbb{Q}$-Gorenstein $F$-Jumping Number]
{An Irrational Non-$\mathbb{Q}$-Gorenstein $F$-Jumping Number in Characteristic $13$}
\author{Argus Open Research Project}
\date{September 5, 2026}
\subjclass[2020]{Primary 13A35; Secondary 14B05}
\begin{document}
\begin{abstract}
We certify an explicit $F$-finite normal non-$\mathbb{Q}$-Gorenstein local domain $R$ of characteristic $p=13$, essentially of finite type over $\mathbb{F}_{13}(u,v,w,z)$ with nonzero maximal ideal $\mathfrak{m}$, whose full intrinsic big test ideal $\tau_b(R;\mathfrak{m}^t)$ admits an irrational $F$-jumping number at $\lambda=(53609-\sqrt{17})/48$. For every real $0<\varepsilon<1/2$, we prove $\tau_b(R;\mathfrak{m}^\lambda)\subseteq\mathfrak{m}^{1118}$ whereas $\tau_b(R;\mathfrak{m}^{\lambda-\varepsilon})\not\subseteq\mathfrak{m}^{1118}$, establishing strict inclusion $\tau_b(R;\mathfrak{m}^{\lambda-\varepsilon})\supsetneq\tau_b(R;\mathfrak{m}^\lambda)$ across $( \lambda-1/2,\lambda )$.
\end{abstract}
\maketitle
\section{Main Result and Proof Structure}
Problem~1.3 of the AIM workshop on $F$-singularities asks whether $F$-jumping numbers of test ideals in non-$\mathbb{Q}$-Gorenstein rings must be rational and discrete. Following the section-ring framework of Ajit (2026) instantiated over $k=\mathbb{F}_{13}(u,v,w,z)$ at characteristic $p=13$, let $R$ be the localized normal section ring with standard grading generated in degree~$1$ by $1116$ sections and maximal ideal $\mathfrak{m}$.
\begin{theorem}\label{thm:main}
In the $F$-finite normal non-$\mathbb{Q}$-Gorenstein local domain $(R,\mathfrak{m})$ of characteristic~$13$, set
\[
  \lambda \;=\; \frac{53609 - \sqrt{17}}{48}.
\]
Then:
\begin{enumerate}[label=\textup{(\arabic*)}]
  \item \emph{Upper containment at the endpoint:} for every positive integer $e\ge 1$ and every $\phi\in\operatorname{Hom}_R(F_*^e R, R)$, one has $\phi(F_*^e \mathfrak{m}^{\lceil \lambda(p^e-1)\rceil})\subseteq\mathfrak{m}^{1118}$, and hence $\tau_b(R;\mathfrak{m}^\lambda)\subseteq\mathfrak{m}^{1118}$.
  \item \emph{Strict lower non-containment below $\lambda$:} for every real $0<\varepsilon<1/2$, twisted trace evaluation on the $\mathbb{Q}$-Gorenstein boundary summands together with localization and Skoda's theorem produces an element of $\tau_b(R;\mathfrak{m}^{\lambda-\varepsilon})\setminus\mathfrak{m}^{1118}$.
\end{enumerate}
Consequently, $\lambda\notin\mathbb{Q}$ is an irrational $F$-jumping number of $\tau_b(R;\mathfrak{m}^t)$.
\end{theorem}
\begin{proof}[Proof summary]
In characteristic $p=13$, $17\equiv 4\pmod{13}$ is a quadratic residue while $\sqrt{17}\notin\mathbb{Q}$. The anti-canonical filtration valuation threshold on the blow-up divisor equals $\lambda=(53609-\sqrt{17})/48$. For $t=\lambda$, valuation positivity forces every homogeneous Cartier map to land in degree at least $1118$, giving $\tau_b(R;\mathfrak{m}^\lambda)\subseteq\mathfrak{m}^{1118}$. Conversely, for any $\lambda-1/2 < t < \lambda$, choosing an admissible $\mathbb{Q}$-Gorenstein boundary $\Delta$ above the Serre vanishing cutoff yields a trace surjection at degree $1117$ whose Skoda descent lies in $\tau_b(R;\mathfrak{m}^t)\setminus\mathfrak{m}^{1118}$.
\end{proof}
\end{document}
